\documentclass[leqno, a4paper, 12pt]{jsarticle}
\usepackage{amssymb}
\usepackage{amsthm}
\usepackage{amsmath}
\usepackage{comment}
\usepackage{url}

\usepackage{enumitem}
%\usepackage{showkeys}


%定理環境 thmはあまり使わずpropをなるべく使う。
%主張は基本的に証明の中で使い、そのため番号を付けない。
%注意環境を付けてみた。
\newtheorem{thm}{定理}
\newtheorem{prop}[thm]{命題}
\newtheorem{cor}[thm]{系}
\newtheorem{lem}[thm]{補題}
\newtheorem{df}[thm]{定義}
\newtheorem{exam}[thm]{例}
\newtheorem{fact}[thm]{事実}
\newtheorem*{claim}{主張}
\newtheorem*{cau}{注意}
\newtheorem*{rmk}{注意}
\renewcommand{\proofname}{証明}
%矢印たち。
%\toは\rightarrowと同じ。\getsは\leftarrowと同じ。同値は\iff
\newcommand{\To}{\Longrightarrow}
\newcommand{\Gets}{\LongLeftarrow}
%黒板太字
\newcommand{\nn}{\mathbb{N}}
\newcommand{\zz}{\mathbb{Z}}
\newcommand{\qq}{\mathbb{Q}}
\newcommand{\rr}{\mathbb{R}}
\newcommand{\cc}{\mathbb{C}}
%無限大
\newcommand{\mgn}{\infty}
%差集合は\setminusで統一する。等号の否定は\neq
%真部分集合は\subsetneqq　偏微分は\partial
%ディスプレイスタイル。\dfracでディスプレイ分数
\newcommand{\dpst}{\displaystyle}

\title{論文メモ
\large{\\--- $n$-cardinality ---}
}
\author{ナンブ　キトラ}
\date{\today}
\begin{document}
\maketitle

本棚を整理したいので，
溜まっている論文のメモを書いて未来への手紙とすることにする．


以下の論文は以下のような空間の構成にまつわる論文を集めたものだと思う：
\begin{thm}
性質$\mathcal{P}$をパラコンパクト，リンデレフ，サブパラコンパクトのどれか一つとする．そして
性質$\mathcal{Q}$は正規か族正規のいずれかとする．
そして$k, m$を$1\le m\le k\le \aleph_{0}$を満たすものとして任意にとる．
すると以下の性質を満たす可分で第一可算な位相空間$X$が存在する．
\begin{enumerate}[label=\textup{(\arabic*)}]
\item 
$X^{n}$が性質$\mathcal{P}$を満たすことと
$n<k$は同値である；

\item 
$X^{n}$が性質$\mathcal{Q}$
を満たすことと$n<m$は同値である．
\end{enumerate}
\end{thm}
この定理の証明は
\cite{MR0561584}
にある．
論文
\cite{MR0540608}
もそういう話．
これにまつわる論文を集めたというか，
大体
Przymusi\'{n}ski
の論文である．

正規性は直積では保たれない．
例えばゾルゲンフライ直線などが
シンプルな例になっている．
パラコンパクト性やリンデレフも
そうである．
上記の定理はそれではそれらの直積で保たれない性質が消えるような臨界的な直積の回数はなんであるかという疑問に答えるものである．いくらでも任意にいじくれるというのが
上記の定理の答えである．

さてこの定理の空間の構成をちゃんと追っかけたわけではないが，空間の構成には$n$-cardinality 
というものを用いる．
これを紹介したい．

集合$X$
に対して自然数$n$回の直積
$X^{n}$を考える．
このとき点
$p=(p_{1}, \dots, p_{n})\in X^{n}$に対して
\[
\hat{p}=\{p_{1}, \dots, p_{n}\}
\]
というふうに$X$の部分集合$\hat{p}$
を対応させる．
例えば$p$が対角線に入っているなら
$\hat{p}$は一点集合になることに注意しよう．
このとき以下が成り立つ．
\begin{prop}[\cite{MR0491191}]
$X$を集合で$n$を自然数とする．
このとき$A$を
$X^{n}$の部分集合として以下の三つの量はどれか一つが無限になるのであれば三つの値は一致する．
\begin{enumerate}[label=\textup{(\arabic*)}]
\item 
$\max\{\, |B|\mid \text{$B$は$A$の部分集合で$B$から異なる2元$p, q$をとると常に$p_{i}\neq q_{i}$になる．}\,\}$; 


\item 
$\max\{\, |B|\mid \text{$B$は$A$の部分集合で$B$から異なる2元$p, q$をとると常に$\hat{p}\cap \hat{q}=\emptyset$になる．} \, \}$; 
\item 
$\max\{\, |Y|\mid 
\text{$Y$は$X$の部分集合で$A\subset \bigcup_{i=1}^{n}(X^{i-1}\times Y\times X^{n-i})$を満たす．}\, \}$

\end{enumerate}
ここで
$|B|$などは$B$の濃度を表す．
\end{prop}

上の命題の二番で定義される量を
$A$の
$n$-cardinality
とよび
$|A|_{n}$と書くことにする．
$A$が$n$-countableであるとは
$|A|_{n}\le \aleph_{0}$
であることとし，
そうでない時には
$n$-uncountable 
と呼ぶ．
つまり
$n$-cardinality
とは，
$A$の点の集合であって座標がダブらないような最大の
$A$の部分集合であるということである．

この概念が先に述べた定理でどのように使われるのかは説明できないけれども，
とにかく大事なのは
$n$-cardinality
に対してもボレル集合の完全集合性質は成り立つということである．
以下が成り立つ．

\begin{thm}[\cite{MR0491191}]
$X$をポーランド空間とし
$n$を自然数とする．
そして$B$を$X^{n}$の
ボレル集合とする．
このとき以下は同値である．
\begin{enumerate}[label=\textup{(\arabic*)}]
\item 
$B$は$n$-uncountable 
である．
\item 
$|B|_{n}=2^{\aleph_{0}}$

\item
$C$をカントール集合とする．
位相的埋め込み写像$h\colon C\to B$
が存在して
$\widehat{h(x)}\cap \widehat{h(y)}=\emptyset$
を任意の$x, y\in C$でみたす．

\item 
$C$をカントール集合とする．
このとき位相的埋め込み写像の族
$\{h_{i}\colon C\to X\}_{i=1}^{n}$が存在し，さらに
$h(x)=(h_{1}(x), \dots, h_{n}(x))$で定義される位相的埋め込み写像
$h\colon C\to X^{n}$に関して
$h(C)\subset B$となる．

\item 
その他いろいろ

\end{enumerate}
\end{thm}
$n$-cardinalityでも
完全集合性質が満たされることは面白いと思う．
また比較的新しい次元論の数学書
\cite{MR3970309}
にも$n$-cardinality
が載っている．

%READ!
%myplain is the same to amsplain style. 
\nocite{*}
\bibliographystyle{myplaindoidoi}
\bibliography{prpr.bib}



\end{document}